%!TEX root = ../../thesis-trier.tex

Here, we investigate the complexity
of the constrained synchronization problem for partially ordered input automata
and when the constraint language is given by a
PDFA with up to two states.
First, in \cref{constrinput:sec:pto:NP}, 
we state a quadratic upper bound for shortest synchronizing words of partially ordered
automata in a given regular constraint language
and conclude containment in \NP\  of the 
corresponding decision problem.
\Cref{constrinput:sec:pto:reduction} states
a technical lemma to simplify monoid constraint languages w.r.t.~partially ordered input automata.
In \cref{constrinput:sec:pto:threecases} we collect
three particular cases of constraint languages giving problems in $\AC^0$, that are \LOGSPACE-complete,
or \NP-hard. Using these lemmas, in \cref{constrinput:sec:two-state-pto},
we give a complexity classification for two-state constraint PDFAs and partially ordered input automata.

\subsection{A Bound and Containment in \texorpdfstring{\NP}{NP}} 
\label{constrinput:sec:pto:NP}

In \cref{constr_input:prop:length_constr_sync_word_in_waa},
we state a quadratic bound on shortest synchronizing
words in a given constraint language.  Then, in \cref{thm:constraint_sync_weakly_acyclic_in_NP}, we conclude 
that the constrained synchronization problem is in \NP\ 
for partially ordered input automata.

\begin{proposition}
	\label{constr_input:prop:length_constr_sync_word_in_waa}
	Let $\semiautomat{A}{Q}{\Sigma}{\delta}$ 
	be a partially ordered automaton with $n$ states and $\mathcal B = (\Sigma, P, \mu, p_0, F)$ be a fixed PDFA.
	Then, a shortest synchronizing word for $\mathcal A$ in $L(\mathcal B)$ 
	has length at most $|P|\cdot \binom{n}{2}$.
\end{proposition}
\begin{proof}
	Let $w = x_1 \cdots x_m \in L(\mathcal B)$
	with $x_i \in \Sigma$ and $S_i = \delta(Q, x_1 \cdots x_i)$  
	for $i \in \{1,\ldots,m\}$ and $S_0 = Q$.
	Suppose $m \ge |P|$ and $S_{i+ |P|} = S_i$
	for some $i \in \{0, 1,\ldots, n\}$.
	Since $\mathcal A$ is partially ordered, 
	no word can induce a non-trivial permutation of a subset of states.
	So, for the word $u = x_{i+1} \cdots x_{i + |P|}$, for each $q \in S_j$,
	we have $\delta(q, u) = q$,
    and then $\applytrans{\delta}{q}{x_j} = q$ for $j\in \{i+1,\ldots,i+|P|\}$.
	Since $u$ has length~$|P|$, it induces a loop in the constraint automaton $\mathcal B$.
	So, we can replace some factor of~$w$
	by a shorter word resulting in a new word $v \in \Sigma^*$ with $|v| < |w|$
	such that
	$ 
	S_{i + |P|} = \delta(Q, x_1\cdots x_i v x_{i + |P| + 1} \cdots x_n)
	$
	and $x_1\cdots x_i v x_{i + |P| + 1} \cdots x_n \in L(\mathcal B)$.

	Now, let $q_1 ,\ldots, q_n$ be a topological sorting of the states of $\mathcal A$ w.r.t.~$\lereach$.
	Suppose $w \in L(\mathcal B)$
	is a shortest synchronizing word for~$\mathcal A$ in $L(\mathcal B)$. 
	Factorize $w$ into blocks of length $|P|$ and a remaining block of length at most $|P|$, i.e.,
	write $w = w_1 \cdots w_m$
	with $|w_i| = |P|$ for $i \in \{1,\ldots,m-1\}$
	and $|w_m| \le |P|$.
	By the arguments from the previous paragraph
	we have $\delta(Q, w_1 \cdots w_{i-1}) \ne \delta(Q, w_1 \cdots w_i)$
	for $i \in \{1,\ldots,m\}$ (set $w_1 \cdots w_0 = \varepsilon$).
	Thus we can define a function $f \colon \{1,\ldots,m\} \to Q \times Q$
	by
	\[
	f(i) = (q_j, \delta(q_j, w_1 \cdots w_i))
	\]
	where $j \in \{1,\ldots,m\}$ is some number
    such that $\delta(q_j, w_1 \cdots w_{i-1}) \ne \delta(q_j, w_1 \cdots w_i)$.
    Since~$\mathcal A$ is partially ordered
	and 
	$\delta(q_j, w_0 \cdots w_{i-1}) \lereach 
	\delta(q_j, w_0 \cdots w_{i})$, we find that $f$ is injective.
	We also have 
	\[
	\applytrans{\delta}{q_j}{w_1 \cdots w_i} \in \{q_{j+1}, q_{j+2}, \ldots, q_n\},
	\]
	which yields that the image of $f$ is contained in the set $\bigcup_{j=1}^n \{q_j\} \times \{q_{j+1},\ldots,q_n\}$
	of size $(n-1) + (n-2) + \cdots + 1 = \binom{n}{2}$.
	This implies $m \le \binom{n}{2}$ and so $|w| \le |P| \cdot \binom{n}{2}$.
	
	%
	% Der Beweis ist auch okay. Vielleicht sogar etwas anschaulicher. Nutzt auch den ersten Paragraphen
	%
%Let $q_1 ,\ldots, q_n$ be a topological sorting of the states of $\mathcal A$.
%We represent the situation after reading a word $u \in \Sigma^*$, i.e., the set $\delta(Q, u)$,
%by a tuple $(i_1, \ldots, i_n) \in \mathbb \{1,\ldots,n\}^n$,
%where $i_j$ is the index of $\delta(q_j, u)$ in the topological sorting, i.e.,
%$\delta(q_j, u) = q_{i_j}$.
%Then, $u \in \Sigma^*$ is synchronizing if and only if
%the corresponding tuple is $(n, \ldots, n)$.
%The starting tuple is $(1, \ldots, n)$.
%For $(i_1, \ldots, i_n), (j_1, \ldots, j_n) \in \{1,\ldots,n\}^n$
%we write $(i_1, \ldots, i_n) < (j_1, \ldots, j_n)$
%if, for all $r \in \{1,\ldots,n\}$, we have $i_r \le j_r$
%and there exists at least one $s \in \{1,\ldots, n \}$
%such that $i_s < j_s$.
%	
%	Now, suppose $w = x_1 \cdots x_m \in L(\mathcal B)$ is a shortest synchronizing word for $\mathcal A$. 
%	By the previous paragraph, we can suppose 
%	$S_{i + |P|} \ne S_i$ for every $i \in \{1,\ldots, n\}$.
%	As~$\mathcal A$ is partially ordered and we can only move forward in the topological sorting,
%	if $\delta(Q, u) \ne \delta(Q, uv)$,
%	then for the tuple $(i_1, \ldots, i_n)$ corresponding
%	to $\delta(Q, u)$ and for the tuple $(j_1, \ldots, j_n)$ for $\delta(Q, uv)$
%	we have $(i_1, \ldots, i_n) < (j_1, \ldots, j_n)$. Note that we have equality
%	if and only if $\delta(Q, u) = \delta(Q, uv)$.
%	As we start with $(1, \ldots, n)$
%	and want to reach $(n, \ldots, n)$, %and every time we have $S_i \ne S_{i+1}$
%	%we increase at least one component of the tuple corresponding to $S_i$,
%	we have to increase at least $n - 1$ times the first entry, $n - 2$ times the second
%	and so on.
%	Now, by the previous reasoning, every~$|P|$ symbols we can suppose we increase
%	some component. Combining these observations yields that a shortest
%	synchronizing word has length at most
%	\[
%	|P| \cdot \left( (n - 1) + (n - 2) + \ldots + 1 \right) = |P| \cdot \binom{n}{2}.
%	\]
%	This finishes the proof.
\end{proof}

By Proposition~\ref{constr_input:prop:length_constr_sync_word_in_waa} we can conclude that
for partially ordered input automata, the constrained synchronization problem is always in $\NP$.

\begin{theorem}
	\label{thm:constraint_sync_weakly_acyclic_in_NP}
	For partially ordered input automata and an arbitrary constraint automaton,
	the constrained synchronization problem	is in $\NP$.
\end{theorem}
\begin{proof}
	By Proposition~\ref{constr_input:prop:length_constr_sync_word_in_waa} we can guess a shortest synchronizing word in $L(\mathcal B)$
	of polynomial length.
	Verifying that such a word is indeed synchronizing can be done in polynomial time as \DFASW\ 
	is in \LOGSPACE by \cref{prelim:thm:finite_list_sync_word}.	
\end{proof}

Lastly, we also give a bound for shortest words in a constraint language mapping
one subset into another subset that, as a special case, also yields a quadratic bound
for shortest synchronizing words.

\begin{proposition}
	\label{compl:prop:ST-pto-bound}	
	Suppose $\automat{B}{P}{\Sigma}{\mu}{p_0}{F}$ is a PDFA.
	Let $\semiautomat{A}{Q}{\Sigma}{\delta}$ be a partially ordered automaton 
	and $S, T \subseteq Q$. 
	Then for a shortest word $u \in L(\mathcal B)$
	such that $\applytrans{\delta}{S}{u} \subseteq T$
	we have $|u| < |P||S|(|Q|-1) + |P|$.	
\end{proposition}
\begin{proof}
	We can suppose $S$ and $T$ are non-empty. 
	Suppose $u \in L(\mathcal B)$
	is a shortest word such that $\applytrans{\delta}{S}{u} \subseteq T$.
	Write $u = u_1 \cdots u_m u_{m+1}$ with $|u_1| = \cdots = |u_{m}| = |P|$
	and $|u_{m+1}| < |P|$.
	Assume there exists $i \in \{1,\ldots,m\}$
	such that
	$\applytrans{\delta}{s}{u_1 \cdots u_{i-1}} \eqreach \applytrans{\delta}{s}{u_1 \cdots u_{i-1} u_{i}}$
	for all $s \in S$.\footnote{We set $u_1 \cdots u_0 = \varepsilon$.}
	Since $\mathcal A$ is partially ordered,
	we find $\applytrans{\delta}{s}{u_1 \cdots u_{i-1}} = \applytrans{\delta}{s}{u_1 \cdots u_{i-1} v}$
	for each prefix $v$ of $u_i$ and each $s \in S$.
	Since $|u_i| = |P|$,
	some state of~$\mathcal B$ is repeated when reading~$u_i$
	and thus a non-empty factor of~$u_i$ can be removed
	to obtain a shorter word $w \in \Sigma^*$
	with $\mu(u_1 \cdots u_{i-1} u_i) = \mu(u_1 \cdots u_{i-1} w)$
	and $\applytrans{\delta}{s}{u_1 \cdots u_{i-1}w} = \applytrans{\delta}{s}{u_1 \cdots u_{i-1} u_i}$
	for each $s \in S$.
	Then $u_1 u_{i-1} w u_{i+1} \cdots u_m u_{m+1} \in L(\mathcal B)$
	and $\applytrans{\delta}{S}{u_1 \cdots u_{i-1} w u_{i+1} \cdots u_m} \subseteq T$,
	contradicting the minimal choice of~$u$.
	So, for each $i \in \{1,\ldots,m\}$ there exists $s \in S$
	such that $\applytrans{\delta}{s}{u_1\cdots u_{i-1} u_{i}} \not\lereach \applytrans{\delta}{s}{u_1\cdots u_{i-1}}$.
	Now, every $s \in S$
	can be mapped to a different state at most~$|Q| - 1$ times until reaching a maximal state.
	Hence, for $s \in S$ 
	and the maximal set of indices $1 \le j_1 < \cdots < j_r \le m$
	such that $\applytrans{\delta}{s}{u_1 \cdots u_{j_l-1}} \ne \applytrans{\delta}{s}{u_1 \cdots u_{j_l}}$
	for $l\in \{1,\ldots,r\}$
	we have $r \le |Q| - 1$.
	Since each $i \in \{1,\ldots,m\}$
	appears as such an index for some state from~$S$,
	we find $m \le |S|(|Q|-1)$
	and thus $|u| < |P| |S| (|Q| - 1) + |P|$.
\end{proof}


\subsection{A Reduction Lemma}
\label{constrinput:sec:pto:reduction}

Here, in \cref{constr_sync:lem:pto_monoid_lang}, we prove that for partially ordered input automata
constraint languages that are monoids can be replaced
by languages of the form $\Gamma^*$ for~$\Gamma \subseteq \Sigma$.

\begin{lemma}
	\label{constr_sync:lem:pto_monoid_lang}
	Suppose $L \subseteq \Sigma^*$ is a language that is a monoid
	and $\Gamma = \{ a \in \Sigma \mid L \cap \Sigma^* a \Sigma^* \ne \emptyset\}$.
	If $\semiautomat{A}{Q}{\Sigma}{\delta}$
	is a partially ordered automaton,
	then~$\mathcal A$ has a synchronizing word in $L$
	if and only if $\mathcal A$ has a synchronizing word in $\Gamma^*$.
\end{lemma}
\begin{proof}	
	First, suppose $w \in \Gamma^*$
	synchronizes $\mathcal A$.
	Let $q_1, \ldots, q_n$ be a topological sorting
	of the states. Then $\applytrans{\delta}{Q}{w} = \{q_n\}$
	and	for each state $q_i$ with $i \in \{1,\ldots,n-1\}$
	there exists a letter $a_i \in \Gamma$
	such that
	\[
	 w \in \Sigma^* a_i \Sigma^* \text{ and } \applytrans{\delta}{q_i}{a_i} \ne q_i.
	\]
	By choice of $\Gamma$, for each $a_i \in \Gamma$
	there exists a word $u_i \in L \cap \Sigma^* a_i \Sigma^*$.
	Then $\applytrans{\delta}{q_i}{u_{i}} \ne q_i$
	and so 	$\applytrans{\delta}{q_i}{u_{i}} \in \{q_{i+1}, \ldots, q_n \}$.
	Hence $\delta(Q, u_1 \cdots u_i) \subseteq \{q_{i+1},\ldots, q_n\}$ for $i \in \{1,\ldots,n-1\}$
	and the word $u = u_1 \cdots u_{n-1}$
	synchronizes $\mathcal A$. As $L$ is a monoid, we have $u \in L$.
	
	Conversely, we have $L \subseteq \Gamma^*$ 
	so that every synchronizing word in $L$ is also a synchronizing word in $\Gamma^*$.
\end{proof}

\begin{remark} \label{constrinput:rem:ptomonoid}
	In the proof of \cref{constr_sync:lem:pto_monoid_lang},
	each letter is replaced by words containing this letter.
	This is done for each state, and we cannot simply
	replace letters by words containing those letters
	and obtain a synchronizing words in a partially ordered automaton,
	but must take the states moved by these letters into account.
	For example, consider the partially ordered automaton 
	in \cref{constrinput:fig:ptoreduction}.
	It has $ab$ as a synchronizing word, but $cacb$
	is not a synchronizing word.
	Here, $a$ moves the states $q_1$ and $q_3$
	and~$b$ moves the state~$q_2$.
	Considering the states $q_1, q_2, q_3$
	in that order, we obtain the synchronizing word $cacbca$
	from the method of proof of \cref{constr_sync:lem:pto_monoid_lang}.	
	
	\begin{figure}[htbp]
		\centering
		\begin{tikzpicture}[node distance = 2cm, on grid, auto, ->]
			\node (q1) [state] {$q_1$};
			\node (q2) [state, right of = q1] {$q_2$};
			\node (q3) [state, right of = q2] {$q_3$};
			\node (q4) [state, right of = q3] {$q_4$};
		
			\path[->] (q1) edge node {$c$} (q2)
			(q2) edge node {$c$} (q3);
			
			\path[->] (q1) edge [bend left] node {$a$} (q4)
			(q3) edge node {$a$} (q4)
			(q2) edge [bend right=50, below] node {$b$} (q4);
		\end{tikzpicture}
		\caption{Partially ordered automaton with synchronizing word $ab$, but $cacb$ is not a synchronizing
		word, see \cref{constrinput:rem:ptomonoid}.}
		\label{constrinput:fig:ptoreduction}
	\end{figure}
\end{remark}

\begin{remark} 	Let $\semiautomat{A}{Q}{\Sigma}{\delta}$ be an $n$-state partially ordered automaton
	and $\Gamma \subseteq \Sigma$ with $w \in \bigcap_{a \in \Gamma} \Sigma^* a \Sigma^*$.
	Then for $u \in \Gamma^*$ and $v \in \Sigma^*$
	we have the following equivalences:	
	\begin{align*}
		\delta(Q, uv) = \{q\}
		& \Leftrightarrow \delta(Q, u) \subseteq \delta^{-1}(q, v) \\
		& \Leftrightarrow \text{$\delta^{-1}(q, v)$ contains all trap states of $\mathcal A_{|\Gamma}$} & \tag{Cor.~\ref{cor:pto-trap-states}} \\
		& \Leftrightarrow \delta(Q, w^{n - 1}) \subseteq \delta^{-1}(q, v) & \tag{Prop.~\ref{prelim:pto:prop:power-of-word-to-trap-state}} \\
		& \Leftrightarrow \delta(Q, w^{n-1}v) = \{q\}.		
	\end{align*}
	This implies that for $L \subseteq \Sigma^*$, 
	$\mathcal A$ has a synchronizing word in $\Gamma^* L$
	if and only if it has one in $w^{n - 1}L$.
\end{remark}

\subsection{Three Special Cases} 
\label{constrinput:sec:pto:threecases}

Here, we give three lemmas stating sufficient
conditions for containment in~$\AC^0$ (\cref{constrinput:pto:lem:Sigma1starSigma2}),
\LOGSPACE-completeness (\cref{constrinput:pto:lem:Sigma1starSigma2bstar}),
and for \NP-hardness (\cref{constrinput:pto:lem:Sigma1starSigma2Sigma3star}).
We use these lemmas in our classification
for two-state constraint PDFAs in \cref{constrinput:sec:two-state-pto}. 

\begin{lemma}\label{constrinput:pto:lem:Sigma1starSigma2}
	Let $\Sigma_1, \Sigma_2 \subseteq \Sigma$. %with $\Sigma_1 \cap \Sigma_2 = \emptyset$.
	Then, for a partially ordered automaton
	the problem to decide whether it has a synchronizing
	word in $\Sigma_1^* \Sigma_2$ is in $\AC^0$.	
\end{lemma}
\begin{proof}
	By \cref{constr_sync:thm:union} and as $\Sigma_1^* \Sigma_2 = \bigcup_{a \in \Sigma_2} \Sigma_1^*a$
	we can assume $\Sigma_2 = \{a\}$.
	Let $\semiautomat{A}{Q}{\Sigma}{\delta}$ be a partially ordered input automaton.	
	By \cref{compl:cor:pto_sync_trap},
	$\mathcal A$ has a synchronizing word in $\Sigma_1^*a$
	if and only if there exists a word in $\Sigma_1^*a$
	mapping all states to a unique trap state $t \in Q$. 
	This is equivalent to the existence of
	a word $u \in \Sigma_1^*$ such that $\delta(Q, u) \subseteq S$ where $S = \delta^{-1}(\{t\}, a)$.	
	We assume some numbering $Q = \{1,\ldots,n\}$ of the states
	and that $\mathcal A$
	is given by a set of transition matrices $(m^{(a)}_{ij})_{i,j \in Q}$
	for each letter $a \in \Sigma$.
	Let $UniqueTrap_j = \left( \bigwedge_{a \in \Sigma} m^{(a)}_{jj} \right)
	\land 
	\left( \bigwedge_{r \in Q \setminus \{j\}} \bigvee_{a \in \Sigma} \neg m_{rr}^{(a)} \right)$
	be an abbreviation for the formula that checks 
	whether $j$ is a unique trap state.
	Having a trap state $t \in Q$ in $\mathcal A$,
	a state $i \in Q$ is in $S$ if and only if the following formula is true:	
	\[
	\bigvee_{j \in Q} \left[ UniqueTrap_j \land m^{(a)}_{ij} \right].
	\]
	So, we can check whether $\mathcal A$ has a unique trap state and
	determine the states in~$S$ by an $\AC^0$ circuit. Then,
	by feeding the input $(\mathcal A_{|\Sigma_1}, S)$
	to \SIS, we obtain an $\leACZEROm$-reduction
	to the problem \SIS on partially ordered automata,
	which is in $\AC^0$ by \cref{compl:thm:sync_into_pto}.
\end{proof}

\begin{lemma} \label{constrinput:pto:lem:Sigma1starSigma2bstar}
	Let $\Sigma_1, \Sigma_2 \subseteq \Sigma$ and $b \in \Sigma$ with $( \Sigma_1 \cup \Sigma_2 ) \setminus \{b\} \ne \emptyset$.
	Then the problem to test whether a given partially ordered
	automaton has a synchronizing word in $\Sigma_1^*\Sigma_2 b^*$
	is \LOGSPACE-complete.
\end{lemma} 
\begin{proof}
	\emph{Containment in \LOGSPACE:}
	By \cref{constr_sync:thm:union} and as $\Sigma_1^* \Sigma_2 b^* = \bigcup_{a \in \Sigma_2} \Sigma_1^*a b^*$
	we can assume $\Sigma_2 = \{a\}$.
	Let $\semiautomat{A}{Q}{\Sigma}{\delta}$ be a partially ordered automaton
	with~$n$ states. By \cref{prelim:classes:eqn:pto} and \cref{prelim:thm:sync_words_is_ideal} 
	a word $uab^i$ with $i \ge 0$ and $u \in \Sigma_1^*$ synchronizes~$\mathcal A$
	if and only if $uab^n$ synchronizes~$\mathcal A$.
	By \cref{compl:cor:pto_sync_trap}, the synchronizing state
	is a unique trap state. 
	So,~$\mathcal A$ has a synchronizing word in $\Sigma_1^* a b^*$
	if and only if there exists a unique trap state $t \in Q$ 
	and a word $u \in \Sigma_1^*$ such that $\delta(Q, u) \subseteq S$
	where $S = \delta^{-1}(\{t\}, ab^n)$.
	By iterating through the states
	we can check whether $\mathcal A$ has a unique trap state $t \in Q$.
	Then, by iterating through the states $q \in Q$ again
	and computing, by using a counter and a memory cell,
	$\delta(q, ab^n)$ 
	we can compute the set $S$.
	Then, solve \SIS with input $(\mathcal A_{|\Sigma_1}, S)$,
	which is in $\AC^0$ by \cref{compl:thm:sync_into_pto}
	as $\mathcal A_{|\Sigma_1}$ is partially ordered.
		
	\medskip 
	\noindent\emph{\LOGSPACE-hardness:} We give an $\leACZEROm$-reduction
	from $\ONEDAG$, which is \LOGSPACE-complete by \THM~\ref{prelim:compl:thm:ONEDAG}.
	Suppose $(G, s, t)$ is a $\ONEDAGST$-instance
	with $G = (V, E)$.
	We can suppose~$s$ has no incoming transitions and~$t$ has no outgoing transitions.
	Pick $u \in \Sigma_{1}^*\Sigma_{2}$ with $|u|_a > 0$ for some $a \in (\Sigma_1 \cup \Sigma_2) \setminus \{b\}$.
		
	\medskip
	\noindent\underline{Construction:}
	Let $\semiautomat{A}{V}{\Sigma}{\delta}$ with 
	\[
	\delta(p, a) = \begin{cases}
		s  & \text{if } p \ne t, \\
		t  & \text{if } p = t.
	\end{cases}
	\qquad 
	\delta(p, b) = \begin{cases} 
		q & \exists (p, q) \in E, \\
		p & \text{otherwise,}  
	\end{cases} 
	\]
	and set $\delta(q, x) = q$ for all $x \in \Sigma \setminus \{a,b\}$.	
	
	\medskip 
	\noindent\underline{Properties:} 
	\begin{enumerate}
		
		\item \label{constrinput:pto:lem:Sigma1starSigma2bstar:1} For each $q \in Q$ we have $\delta(q, u) = \delta(q, a)$.
		
		This follows since $|u|_b = 0$, only the letters $a$ and $b$
		move any state, and $\delta(q, a) = \delta(q, aa)$ for each $q \in Q$.		
		
		\item \label{property:areset} For $v \in \Sigma^*$ and $|\delta(Q, v)| \ge 2$
		we have $\delta(Q, va) = \{s,t\}$.
		
		As $t$ is a trap state, $t \in \delta(Q, v)$.
		For $q \in \delta(Q, v)$ with $q \ne t$	we have $\delta(q, a) = \{s\}$.	This yields the claim.		
		
		\item \label{property:path}There exists a word $v \in \Sigma^*$
		with $\delta(\{s,t\}, v) = \{t\}$ if and only if there exists a path of length $n$ from $s$ to $t$
		in $G$ and in this case $\delta(\{s,t\}, b^n) = \{t\}$.
		
		If there exists a path of length $n$ from $s$ to $t$ in $G$,
		then $\delta(\{s,t\}, b^n) = \{t\}$.
		Conversely, suppose $v \in \Sigma^*$ is a shortest word with $\delta(\{s,t\}, v) = \{t\}$.
		By Property~\ref{property:areset} 
		we can assume $|v|_a = 0$. As the other letters only induce
		self-loop, we have $v \in b^*$.
		Now, by following the transitions along $b$ we obtain a path in $G$ from $s$ to $t$, i.e.,
		the states $s, \delta(s, b), \delta(s, bb), \ldots, \delta(s, b^{|v|-1}), t$
		give a path in $G$.		
	\end{enumerate}
	
	\noindent\underline{Verification:} Suppose there exists
	a path of length $n$ from $s$ to $t$ in $G$.
	By Properties~\ref{constrinput:pto:lem:Sigma1starSigma2bstar:1}, \ref{property:areset},
	and~\ref{property:path}, $\delta(Q, ub^n) = \{t\}$.	
	Conversely, suppose $v \in \Sigma_{1}^* \Sigma_{2} b^*$ 
	synchronizes $\mathcal A$. 
	Then $\delta(Q, v) = \{t\}$.
	In particular $\delta(\{s,t\}, v) = \{t\}$. Thus, by Property~\ref{property:path},
	there exists a path from $s$ to $t$ in $G$. 
\end{proof}

\begin{lemma}\label{constrinput:pto:lem:Sigma1starSigma2Sigma3star}
	Let $\Sigma_1, \Sigma_2, \Sigma_3 \subseteq \Sigma$
	be such that $\Sigma_1 \cap \Sigma_2 = \emptyset$ and
	there exist distinct letters $b,c \in \Sigma_{3}$
	and $a \in (\Sigma_1 \cup \Sigma_2) \setminus \Sigma_3$.
	Then the problem to decide whether a given 
	partially ordered automaton has a synchronizing word in $\Sigma_1^*\Sigma_2\Sigma_3^*$
	is \NP-hard.
\end{lemma}
\begin{proof}	
	We give a reduction from the \NP-complete problem~{\sc SAT} (Theorem~\ref{prelim:thm:complexities})
	that is a modification of a reduction from Eppstein~\cite{Eppstein90} 
	and Rystsov~\cite{Rys80}
	used to show \NP-hardness of \DFASWS. A similar modification was used in~\cite[Thm.~4]{Ryzhikov19a}
	to show that \SFS is \NP-hard for partially ordered 
	input automata.	
		
	Let $(X,C)$ be a {\sc SAT}-instance with variables $X = \{x_1, \ldots, x_n\}$
	and clauses $C = \{c_1, \ldots, c_m\}$.
	
\medskip 
\noindent\underline{Construction:} We construct a partially ordered automaton $\semiautomat{A}{Q}{\Sigma}{\delta}$
with $Q = \{ q_{i,j} \mid 1 \le i \le m, 0 \le j \le n + 1 \}\cup \{q_f\}$
where $q_f$ is a trap state and $\delta \colon Q \times \Sigma \to Q$  is given, 
for $i \in \{1,\ldots,m\}$, $j \in \{0,1,\ldots,n+1\}$ and $x \in \Sigma$, by 
\[
\applytrans{\delta}{q_{i,j}}{x} = \begin{cases}
q_{i, j + 1} & \mbox{if } j \in \{1,\ldots, n\} \land ( ( x = b \land \neg x_j \notin c_i ) \lor ( x= c \land  x_j \notin c_i ) ), \\
q_f          & \mbox{if } j \in \{1,\ldots, n\} \land ( ( x = b \land \neg x_j \in c_i ) \lor ( x = c \land    x_j \in c_i ) ), \\ 
%q_f          & \mbox{if } j \in \{1,\ldots, n\} \land x = c \land \overline x_j \in c_i \\
q_{i,1}      & \mbox{if } j = 0 \land x \in \Sigma_2, \\ 
q_f          & \mbox{if } j \in \{1,2,\ldots,n+1\} \land x = a, \\
q_{i,j}      & \mbox{otherwise},
\end{cases}
\]
and $\applytrans{\delta}{q_f}{x} = x$. See Figure~\ref{fig:extended_Eppstein} for an example.

	\begin{figure}[ht]
	\centering
		\begin{tikzpicture}[transform shape,scale=0.8,shorten >=1pt,node distance=2.1cm and 3.5cm,on grid,auto, state/.style={circle, draw, blur shadow,fill=white,minimum size=1.3cm, fill=black!5, thick}] % alternative: adding "minimum size=1.5" as option to each node
			
			% formula above the automaton
		\node (clause_1) {\large$(x_1 \vee x_2)$};
		\node[right=of clause_1] (clause_2) {\large $(x_1 \vee \neg x_2)$};
		\node[right=of clause_2] (clause_dots) {$\wedge$};%{\small \dots};
		\node[right=of clause_dots] (clause_{m-1}) {\large $(\neg x_1)$}; % \vee x_3)$};
		\node[right=of clause_{m-1}] (clause_m) {\large $(x_1 \vee x_2)$};
		% "Logical Or"-labeling-nodes of the formula over the automaton
		\node (label_wedge_1) at ($(clause_1)!0.48!(clause_2)$) {$\wedge$}; % measurement of the length seems so start at the outermost points of the nodes, not from the mids => check it manually, especially if the node diameter changes
		\node (label_wedge_2) at ($(clause_2)!0.52!(clause_dots)$) {};%{$\wedge$};
		\node (label_wedge_{m-2}) at ($(clause_dots)!0.46!(clause_{m-1})$) {};%{$\wedge$};
		\node (label_wedge_{m-1}) at ($(clause_{m-1})!0.535!(clause_m)$) {$\wedge$};
		
	
		
		% column 3ff (node order has to be this way, because only known nodes can be referenced correctly - otherwise the positioning breaks)
		%\node[below=of q_s] (q_{3ff,1}) {$\dots$};
		\node[below=of clause_dots] (q_{3ff,0}) {};%{$\dots$};
		\node[below=of q_{3ff,0}] (q_{3ff,1}) {};%{$\dots$};
		\node[below=of q_{3ff,1}] (q_{3ff,2}) {};%{$\dots$};
		\node[below=of q_{3ff,2}] (q_{3ff,3}) {};% {$\dots$};		
		
		% column 2
		\node[state, left=of q_{3ff,0}, yshift=0.6cm] (q_{2,0}) {$q_{2,0}$};
		\node[state, below=of q_{2,0}] (q_{2,1}) {$q_{2,1}$};
		\node[state, below=of q_{2,1}] (q_{2,2}) {$q_{2,2}$};
		\node[state, below=of q_{2,2}] (q_{2,3}) {$q_{2,3}$};		
		
		% column 1
		\node[state, left=of q_{2,0}] (q_{1,0}) {$q_{1,0}$};
		\node[state, left=of q_{2,1}] (q_{1,1}) {$q_{1,1}$};
		\node[state, below=of q_{1,1}] (q_{1,2}) {$q_{1,2}$};
		\node[state, below=of q_{1,2}] (q_{1,3}) {$q_{1,3}$};		
		
		% column n-1
		\node[state, right=of q_{3ff,0}, yshift=0.6cm] (q_{m-1,0}) {$q_{3,1}$};
		\node[state, below=of q_{m-1,0}] (q_{m-1,1}) {$q_{3,1}$};
		\node[state, below=of q_{m-1,1}] (q_{m-1,2}) {$q_{3,2}$};
		\node[state, below=of q_{m-1,2}] (q_{m-1,3}) {$q_{3,3}$};		
		
		% column n
		\node[state, right=of q_{m-1,0}] (q_{m,0}) {$q_{4,0}$};
		\node[state, below=of q_{m,0}] (q_{m,1}) {$q_{4,1}$};
		\node[state, below=of q_{m,1}] (q_{m,2}) {$q_{4,2}$};
		\node[state, below=of q_{m,2}] (q_{m,3}) {$q_{4,3}$};		
		
		% q_t 
		%\node[state, below=of q_{3ff,n+1}](q_t) {$q_t$};
		\node[state, below=of q_{3ff,3}, yshift=-0.3cm](q_t) {$q_f$};		
		
		\path[use as bounding box] (0, 0) rectangle (0, 0); % (remove the border that arises through the invisible nodes, which are needed for the curves)	
		
		
		% edges of column 2 (come first, so that the edges are drawn below the loops of column 1)
		\path[->, every loop/.style={looseness=3}]
		(q_{2,0}) edge node {$x \in \Sigma_2$}  (q_{2,1})
		(q_{1,0}) edge node {$x \in \Sigma_2$} (q_{1,1})
		(q_{m-1,0}) edge node {$x \in \Sigma_2$}  (q_{m-1,1})
		(q_{m,0}) edge node {$x \in \Sigma_2$}  (q_{m,1})
		;
		\path[->, every loop/.style={looseness=3}]
		(q_{2,1})
		edge[out=0, in=190, controls=+(180:1.2cm) and +(205:7.5cm)] (q_t)
		edge[densely dashed] (q_{2,2})
		%edge[-,draw=white,line width=7pt,loop right] node[line width = 3pt,draw=white,fill=white, text=white, ultra thin] {$\sigma_{k'}$} ()
		%edge[out=25,in=335,looseness=5,very thick,dotted] (q_{2,1})
		edge[out=0, in=140, very thick, dotted] (q_t)
		(q_{2,2})
		edge[densely dashed, out=0, in=180, controls=+(180:1.2cm) and +(205:6.85cm)] (q_t)
		edge (q_{2,3})
		%edge[-,draw=white,line width=7pt,loop right] node[line width = 3pt,draw=white,fill=white, text=white, ultra thin] {$\sigma_{k'}$} ()
		%edge[loop right,gray,draw=black!25] node[] {$\sigma_{k'}$} ()
		edge[very thick,dotted,out=0,in=140] (q_t)
		(q_{2,3}) edge[very thick,dotted,out=330,in=200] (q_t)
		;
		
		% edges of column 1 (loops go over the edges of column 2)
		\path[->, every loop/.style={looseness=3}] % "text=gray"/usage of different grayscales to emphasize the text optically
		(q_{1,1})
		edge[out=0, in=240, controls=+(180:2cm) and +(200:12.5cm)] (q_t) % this notation works too
		edge[densely dashed] (q_{1,2})
		% 	edge[-,draw=white,line width=7pt,loop right] node[line width = 3pt,draw=white,fill=white, text=white, ultra thin] {$\sigma_{k'}$} () % the edge is drawn in white first to get a border around it, ...
		%edge[out=25,in=335,looseness=5,very thick,dotted] (q_{1,1}) % ... then the edge is drawn
		edge [very thick,dotted, out=295, in=200] (q_t)
		(q_{1,2})
		edge[out=0, in=230, controls=+(180:2cm) and +(200:11.25cm)] (q_t)
		edge[densely dashed] (q_{1,3})
		% 	edge[-,draw=white,line width=7pt,loop right] node[line width = 3pt,draw=white,fill=white, text=white, ultra thin] {$\sigma_{k'}$} ()
		%	edge[loop right,gray,draw=black!25] node[] {$\sigma_{k'}$} ()
		%edge [very thick,dotted, out=0, in=240, controls=+(180:2cm) and +(200:12.5cm)] (q_t)
		edge [very thick,dotted, out=300, in=200] (q_t)
		(q_{1,3}) edge[very thick,dotted,out=330,in=200] (q_t)
		;
		
		
		% edges of column n-1
		\path[->, every loop/.style={looseness=3}]
		(q_{m-1,1})
		edge[densely dashed, out=180, in=350, controls=+(0:1.2cm) and +(335:7.5cm)] (q_t)
		edge (q_{m-1,2})
		edge[-,draw=white,line width=7pt,loop left] node[line width = 3pt,draw=white,fill=white, text=white, ultra thin] {$\sigma_{k'}$} ()
		%edge[out=155,in=205,looseness=5,very thick,dotted] (q_{m-1,1})
		edge[very thick,dotted,out=180,in=35] (q_t)
		(q_{m-1,2})
		edge[bend right=10] (q_{m-1,3})
		edge[densely dashed, bend left=10] (q_{m-1,3})
		% 	edge[-,draw=white,line width=7pt,loop left] node[line width = 3pt,draw=white,fill=white, text=white, ultra thin] {$\sigma_{k'}$} ()
		% 	edge[loop left,gray,draw=black!25] node[] {$\sigma_{k'}$} ()
		edge[very thick,dotted,out=180,in=35] (q_t)
		(q_{m-1,3}) edge[very thick,dotted,out=180,in=35] (q_t)
		;
		
		% edges of column n
		\path[->, every loop/.style={looseness=3}]		
		(q_{m,1})
		edge[out=180,in=300, controls=+(0:2cm) and +(340:12.5cm)] (q_t)
		edge[densely dashed] (q_{m,2})
		edge[-,draw=white,line width=7pt,loop left] node[line width = 3pt,draw=white,fill=white, text=white, ultra thin] {$\sigma_{k'}$} ()
		%edge[out=155,in=205,looseness=5,very thick,dotted]  (q_{m,1})
		edge[very thick,dotted,out=245,in=340] (q_t)
		(q_{m,2})
		edge[out=180,in=310, controls=+(0:2cm) and +(340:11.25cm)] (q_t)
		edge[densely dashed] (q_{m,3})
		%edge[-,draw=white,line width=7pt,loop left] node[line width = 3pt,draw=white,fill=white, text=white, ultra thin] {$\sigma_{k'}$} ()
		%edge[loop left,gray,draw=black!25] node[] {$\sigma_{k'}$} ()
		edge[very thick,dotted,out=240,in=340] (q_t)
		(q_{m,3}) edge[very thick,dotted,out=210,in=340] (q_t)
		;		
		
	\end{tikzpicture}
\vspace{11pt} 
\caption[Example for the reduction from the proof of \cref{constr_sync:thm:input_pto}]{Reduction from the proof of \cref{constr_sync:thm:input_pto}
	for the formula
	$(x_1 \lor x_2) \land (x_1 \lor \neg x_2)
	\land (\neg x_1) \land (x_1 \lor x_2)$.
	Transitions %at states $q_{i,j}$ with $j \ge 1$ 
	labeled by $a$, $b$, and~$c$
	are drawn with dotted, dashed, and solid lines, respectively. Self-loops omitted.}
\label{fig:extended_Eppstein}	
\end{figure}



\medskip 
\noindent\underline{Properties:} Properties~\ref{constr_sync:item:pto_qi0}, 
\ref{constrsync:item:pto:proj}, \ref{constrsync:subclasses:pto:item:monotone}
and~\ref{constr_sync:item:pto_assignment} 
follow readily. For the remaining two properties we give the proofs.


\begin{enumerate}
\item \label{constr_sync:item:pto_qi0}
$\applytrans{\delta}{q_{i,0}}{x} \ne q_{i,0}$ if and only if $x \in \Sigma_2$.

\item \label{constrsync:item:pto:proj} For $u \in (\Sigma \setminus \{a\})^*$ 
and $i \in \{1,\ldots,m\}$ 
we have $\applytrans{\delta}{q_{i,1}}{u} = \applytrans{\delta}{q_{i,1}}{\pi_{\{b,c\}}(u)}$.

\item \label{constrsync:subclasses:pto:item:monotone} Let $i \in \{1,\ldots,m\}$ and $j,j' \in \{0,1,\ldots,n+1\}$.
If $\applytrans{\delta}{q_{i,j}}{u} = q_{i,j'}$ for some $u \in \Sigma^*$,
then $j \le j'$.

\item \label{constr_sync:item:pto_assignment} 
Let $i \in \{1,\ldots,m\}$, $v_j \in \{b,c\}$ for $j \in \{1,\ldots,n\}$ and $w \in \Sigma^*$.
Then
\begin{align*}
	& \applytrans{\delta}{q_{i,1}}{v_1 \cdots v_n w} = q_f \\
	& \Leftrightarrow 
	\exists j \in \{1,\ldots,n\} \colon (v_j = b \land x_j \in c_i) \lor ( v_j = c \land \neg x_j \in c_i) \\
	& \Leftrightarrow  \text{the assignment } x_1, \ldots, x_n \text{ satisfies the clause $c_i$, }
\end{align*}
where
\begin{equation}\label{constrsync:input:pto:eqn:assignment}
	x_j = \begin{cases}
		0 & \text{if } v_j = b, \\
		1 & \text{if } v_j = c
	\end{cases}
\end{equation}
for $j \in \{1,\ldots,n\}$.

\item \label{constr_sync:item:pto_assignment2} Let $i \in \{1,\ldots,m\}$ and $v \in \{b,c\}^*$.
Then $\delta(q_{i,1}, v) = q_f$
implies the existence of a satisfying assignment for the clause $c_i$.

Let $w \in \pref{vb^*} \cap \Sigma^n$.
Then $\delta(q_{i,1}, w) = q_f$,
which, for $|v| > n$, follows
as otherwise the word $v$ would end up in $q_{i,n+1}$,
and, for $|v| \le n$, as $q_f$ is a trap state.
Property~\ref{constr_sync:item:pto_assignment} 
then yields the claim.

\item \label{constrsync:subclasses:pto:syncword} $\delta(Q, u) = \{q_f\}$ for $u \in \Sigma^*$ if and only if $u \in \Sigma^* a \Sigma^* \cap (\Sigma \setminus \Sigma_2)^*\Sigma_2 w$
 for some $w \in \Sigma^*$
 with  $\delta(q_{i,1}, w) = q_f$ for all $i \in \{1,\ldots,m\}$.
  
 First, assume $u \in \Sigma^* a \Sigma^* \cap (\Sigma \setminus \Sigma_2)^*\Sigma_2 w$
 with  $\delta(q_{i,1}, w) = q_f$ for all $i \in \{1,\ldots,m\}$.
 Set $S = \{ q_{i,j} \mid i \in \{1,\ldots,m\}, j \in \{1,\ldots,n+1\}\}$.
 By Property~\ref{constrsync:subclasses:pto:item:monotone}, $\delta(S, u) \subseteq S$.
 As $\delta(S, a) = \{q_f\}$ and $u \in \Sigma^*a\Sigma^*$,
 we find $\delta(S, u) = \{q_f\}$.
 As precisely letters from $\Sigma_2$ map $q_{i,0}$ to $q_{i,1}$ and all other letters fix $q_{i,0}$ 
 we get $\delta(q_{i,0}, u) = \delta(q_{i,1}, w) = q_f$
 for all $i \in \{1,\ldots,m\}$.
 Hence, $\delta(Q, u) = \{q_f\}$.
 Conversely, suppose $\delta(Q, u) = \{q_f\}$.
 As the states $q_{i,n+1}$ for $i \in \{1,\ldots,m\}$
 must be mapped to $q_f$ as well, we have $u \in \Sigma^* a \Sigma^*$.  
 By Property~\ref{constr_sync:item:pto_qi0} 
 we get $u = vxw$ for some $v \in (\Sigma\setminus\Sigma_2)^*$, $x \in \Sigma_2$
 and $w \in \Sigma^*$.
 Then $\applytrans{\delta}{q_{i,0}}{vx} = q_{i,1}$
 for all $i \in \{1,\ldots,m\}$.
 So, $\delta(q_{i,1}, w) = q_f$ for $i \in \{1,\ldots,m\}$.
\end{enumerate}

\noindent\underline{Verification:} Let $v \in \Sigma_1^*$, $x \in \Sigma_2$ and $w \in \Sigma_3^*$
and set $S_1 = \{ q_{i,1} \mid i \in \{1,\ldots,m\}\}$.
Then we have the following equivalences:
\begin{align*}
& |\delta(Q, vxw)| = 1 \\
& \Leftrightarrow  \delta(Q, vxw) = \{q_f\} & \tag{$q_f$ is a trap state} \\
& \Leftrightarrow  \delta(S_1, w) = \{q_f\} \land vxw \in \Sigma^* a \Sigma^* & \tag{Property~\ref{constrsync:subclasses:pto:syncword} \& $v \in (\Sigma\setminus\Sigma_2)^*, x \in \Sigma_2$} \\
& \Leftrightarrow  \delta(S_1, w) = \{q_f\} \land vx \in \Sigma^* a \Sigma^* & \tag{$a \notin \Sigma_3$} \\
& \Leftrightarrow  \delta(S_1, \pi_{\{b,c\}}(w)) = \{q_f\} \land vx \in \Sigma^* a \Sigma^*. & \tag{Property~\ref{constrsync:item:pto:proj} \& $a \notin \Sigma_3$} 
\end{align*}  

By Property~\ref{constr_sync:item:pto_assignment2},
$\delta(S_1, \pi_{\{b,c\}}(w)) = \{q_f\}$ for some  $w \in \Sigma_3^*$
if and only if there exists a satisfying assignment for each clause $c_i$, $i \in \{1,\ldots,m\}$. 
So, $|\delta(Q, vxw)| = 1$ with $v \in \Sigma_1^*$, $x \in \Sigma_2$,
and $w \in \Sigma_3^*$ implies that the given {\sc SAT}-instance
is solvable, i.e., a \rYES-instance.
Conversely, suppose there exists a satisfying assignment for each clause.
Then, by Property~\ref{constr_sync:item:pto_assignment},
$\delta(S_1, w) = \{q_f\}$
for some $w \in \{b,c\}^* \subseteq \Sigma_3^*$.
As $a \in \Sigma_1\cup \Sigma_2$ we find $v \in \Sigma_1^*$
and $x \in \Sigma_2$ such that $vx \in \Sigma^* a \Sigma^*$.
Then $|\delta(Q, vxw)| = 1$ with $vxw \in \Sigma_1^*\Sigma_{2}\Sigma_3^*$.
\end{proof}

\subsection{Classification for Two-State Constraint PDFAs}
\label{constrinput:sec:two-state-pto}

Here, in \cref{constr_sync:thm:input_pto}, we classify the complexity for constraint automata
with two states. We find that some constraint languages
giving \PSPACE-hard problems in the general case, as for example
$
(a+b)^*c a^*$ or  
$(a+b)^*cc^*$,
yield \LOGSPACE-complete problems or, like $(a+b)^*c$,
problems in $\AC^0$ when restricted to partially ordered input automata.
But we also find problems that are \PSPACE-hard in the general
case where the complexity only drops to \NP-hardness,
as for example for the constraint language $a(b+c)^*$.
With the assumption from~\cref{constrsync:rem:alphabet}, for one-state automata
the constraint language is either empty or equals $\Sigma^*$.
So, in this case, by \cref{compl:thm:pto_sync_in_AC0}, 
for partially ordered input automata the constrained synchronization
problem is in $\AC^0$.

For a two-state constraint PDFA, we use the notational conventions
as introduced in Section~\ref{sec:constr_sync:small_aut} and the assumption
that the constraint PDFA is trim.

\begin{theorem} 
\label{constr_sync:thm:input_pto}    
    Let $\automat{B}{P}{\Sigma}{\mu}{p_0}{F}$ be a two-state constraint automaton that is trim
    with $P = \{1,2\}, p_0 = 1, F = \{2\}$, and $\Sigma_{i,j} = \{ x \in \Sigma \mid \mu(i, x) = j \}$
    for $i,j \in P$.
    Then the problem to decide whether a given partially ordered automaton $\mathcal A$
    has a synchronizing word in $L(\mathcal B)$
    is
	\begin{enumerate} 
    \item \NP-complete for $\Sigma_{2,1} = \emptyset$,
    $ % beispiele wo einfacher als allgemeiner fall
     ( \Sigma_{1,1} \cup \Sigma_{1,2} ) \setminus \Sigma_{2,2} \ne \emptyset$ and $
       |\Sigma_{2,2}| \ge 2,   
    $
     
     \item \LOGSPACE-complete for $\Sigma_{2,1} = \emptyset$, $(\Sigma_{1,1} \cup \Sigma_{1,2}) \setminus \Sigma_{2,2} \ne \emptyset$ and $|\Sigma_{2,2}| = 1$,  
     
    \item in $\AC^0$ in the remaining cases, which are covered by the following conditions:
    \begin{enumerate}
          
    \item $\mathcal B$ is strongly connected (equivalently $\Sigma_{2,1} \ne \emptyset$), or
    
    \item $\Sigma_{2,2} = \emptyset$, or
        
    \item 	$\Sigma_{1,1} \cup \Sigma_{1,2} \subseteq \Sigma_{2,2}$.
    \end{enumerate}     
	\end{enumerate} 
\end{theorem}
\begin{proof} 
	As $\mathcal B$ is trim we have $\Sigma_{1,2} \ne \emptyset$.
	First, suppose $\Sigma_{2,1} \ne \emptyset$.
	Then $\mathcal B$
	is strongly connected.
	So $\fact{L(\mathcal B)} = \fact{L( \mathcal B_{p_0,\{p_0\}} )}$.
	By \cref{constr_sync:thm:regular_prefix_codes_and_automata}, 
	$L( \mathcal B_{p_0,\{p_0\}} )$ is a monoid.
	By \cref{constr_sync:cor:factors_sets_included,constr_sync:lem:pto_monoid_lang},
	the problem is equivalent 
    to finding a synchronizing word in $\Gamma^*$ for some $\Gamma \subseteq \Sigma$,
    which is in $\AC^0$ by \cref{compl:thm:pto_sync_in_AC0}
    as we can assume $\Gamma = \Sigma$ (see \cref{constr_sync:sec:assumptions}).
	Next, suppose $\Sigma_{2,1} = \emptyset$.
	Then $L(\mathcal B) = \Sigma_{1,1}^* \Sigma_{1,2} \Sigma_{2,2}^*$.
	For $\Sigma_{2,2} = \emptyset$,
	the problem is in $\AC^0$ by \cref{constrinput:pto:lem:Sigma1starSigma2}.
	If $\Sigma_{1,1} \cup \Sigma_{1,2} \subseteq \Sigma_{2,2}$ (and $\Sigma_{2,1}= \emptyset$),
	then $\fact{L(\mathcal B)} = \fact{\Sigma_{2,2}^*}$
	and containment in $\AC^0$ follows from \cref{constr_sync:cor:factors_sets_included} 
	and \cref{compl:thm:pto_sync_in_AC0}
	by considering $\mathcal A_{|\Sigma_{2,2}}$.
	Now, suppose $(\Sigma_{1,1} \cup \Sigma_{1,2}) \setminus \Sigma_{2,2} \ne \emptyset$
	and $\Sigma_{2,2} \ne \emptyset$.
	For $|\Sigma_{2,2}| = 1$,
	\LOGSPACE-completeness follows from \cref{constrinput:pto:lem:Sigma1starSigma2bstar}.	
	For $|\Sigma_{2,2}| \ge 2$,
	\NP-hardness follows from \cref{constrinput:pto:lem:Sigma1starSigma2Sigma3star}.
    By \cref{thm:constraint_sync_weakly_acyclic_in_NP} 
    we can decide in \NP\  whether~$\mathcal A$ has a synchronizing word in 
    $
    L(\mathcal B) = \Sigma_{1,1}^* \Sigma_{1,2} \Sigma_{2,2}^*.
    $	
\end{proof}

Comparing \cref{constr_sync:thm:input_pto} and \cref{thm:two-state-classification}, 
$( \Sigma_{1,1} \cup \Sigma_{1,2} ) \setminus \Sigma_{2,2} \ne \emptyset$
with $|\Sigma_{1,1}| \ge 2$ and $|\Sigma_{2,2}| \le 1$
implies intractability in the general case, whereas
for partially ordered input automata
we obtain a tractable problem. However, for 
$( \Sigma_{1,1} \cup \Sigma_{1,2} ) \setminus \Sigma_{2,2} \ne \emptyset$
and $|\Sigma_{2,2}| \ge 2$,
the constrained synchronization problem is intractable in the general case
and for partially ordered input automata, where we have \PSPACE-completeness
in the former case and \NP-completeness in the latter case.